\section{Introduction}
\label{sec:intro}

Biomedical image segmentation constitutes a fundamental cornerstone of computer-assisted intervention, pathological screening, and automated diagnosis~\cite{litjens2017survey, ronneberger2015unet}. In high-stakes clinical tasks---ranging from early endoscopic detection of precancerous colorectal polyps~\cite{siegel2023colorectal, jha2020kvasir, bernal2017comparative} and dermatological melanoma excision~\cite{codella2019skin} to retinal vascular micro-aneurysm grading~\cite{hoover2000locating, staal2004ridge} and cellular ultrastructure tracing~\cite{cardona2010integrated, caicedo2019nucleus}---accurate delineation of target boundaries is paramount. Automated systems must reliably discern ambiguous tissue interfaces, handle extreme shape heterogeneity, and suppress deceptive background artifacts under severe class imbalance.

To address complex and ill-defined boundaries, deep architectures have increasingly incorporated \textbf{recurrent feedback mechanisms}~\cite{liang2015recurrent, alom2018recurrent, pinheiro2014recurrent, zamir2017feedback}. The foundational intuition of recurrent feedback is inspired by the human visual cortex: biological vision relies heavily on top-down iterative processing, where higher-level semantic interpretations recursively constrain and refine lower-level sensory feature representations~\cite{gilbert2013top, kietzmann2019recurrence}. In medical vision architectures, past prediction masks ($m^{(t-1)}$) generated at feedback step $t-1$ are recursively re-injected into intermediate encoder-decoder feature maps via spatial gating or feature fusion modules. In theory, this recurrent refinement process should iteratively suppress false positives, sharpen ambiguous object contours, and recover subtle morphological details.

\subsection{The Empirical Dilemma: The Feedback Trap and Representation Collapse}
Despite this compelling theoretical premise, practical deployment of recurrent feedback architectures in medical vision reveals a severe, paradoxical pathology: \textbf{chronic false-positive over-segmentation and optimization instability}. In clinical cohorts, qualitative and quantitative audits demonstrate that unconstrained feedback networks frequently exhibit severe prediction bleed, projecting masks tens of pixels beyond genuine anatomical boundaries into healthy parenchyma and background tissue. 

More critically, while modern loss engineering (e.g., asymmetric Tversky loss~\cite{salehi2017tversky} or boundary-weighted losses~\cite{kervadec2019boundary}) reliably curbs false positives in feedforward architectures, its optimization efficacy is completely neutralized once recurrent feedback is engaged. The network becomes trapped in an echo chamber of its own past predictions. 

Through rigorous representational and gradient audits, we identify the mathematical root cause of this failure, which we term the \textbf{Feedback Trap}. The core failure stems from a structural design flaw in conventional recurrent feedback: \textbf{forcing historical predictions to serve simultaneously as forward spatial guidance and as an unconstrained backward optimization path}. In conventional feedback designs, the spatial gating operator typically combines internal feature attention ($\text{fmask}$) with downsampled recurrent feedback ($m_{\text{fg}}$) through hard binary gating:
\begin{equation}
  \label{eq:hard_gate_intro}
  \text{keep}_{\text{hard}} = \max\left(\mathbb{I}(\text{fmask} > 0.5),\; m_{\text{fg}}\right).
\end{equation}
This formulation triggers two structural failures:
\begin{enumerate}
  \item \textbf{Gradient Paralysis of Internal Attention:} Because the indicator function $\mathbb{I}(\cdot > 0.5)$ has a derivative of zero almost everywhere, $\frac{\partial \text{keep}}{\partial \text{fmask}} \equiv 0$. Consequently, the convolutional sub-networks tasked with extracting attention receive zero direct gradient signal from the downstream segmentation loss.
  \item \textbf{Toxic Recursive Backpropagation and Error Amplification:} Concurrently, non-zero gradients backpropagate unconstrained through the recurrent mask branch ($m_{\text{fg}}$). When the network generates an initial false positive, that erroneous prediction is re-injected as a spatial prior in subsequent iterations. Because the backward pass calculates gradients through this recursive dependency graph, early encoder representations experience massive Batch Normalization drift, feature variance collapses, and the model optimizes itself to justify its own prior hallucinations rather than ground-truth anatomy.
\end{enumerate}

Furthermore, empirical observations indicate that without proper stabilization, conventional feedback loops either stagnate entirely across passes (producing zero iterative refinement) or degenerate into runaway over-segmentation when trained over extended epochs.

\subsection{The Solution: Gradient Decoupling and GDF-Net}
To resolve this fundamental pathology, we formulate the \textbf{Feedback Firewall} principle:
\begin{center}
  \textit{Forward spatial guidance must be mathematically decoupled from backward gradient propagation.}
\end{center}
By enforcing a strict stop-gradient operator on recurrent predictions ($\text{detach}(m^{(t-1)})$), historical predictions act strictly as clean forward spatial guidance, completely eliminating recursive error amplification in the backward pass.

Building upon this principle, we introduce \textbf{GDF-Net} (Gradient-Decoupled Feedback Network), an enhanced architecture that unifies three core innovations:
\begin{enumerate}
  \item \textbf{Decoupled Spatial Guidance Stem:} A 4-channel input convolutional stem that ingests both the original image and a detached prediction prior modulated by a learned spatial gate ($\sigma(\mathcal{W}_{\text{gate}} \cdot \text{detach}(m^{(t-1)}))$).
  \item \textbf{Multi-Scale Gradient-Decoupled Feedback Decoder:} Instead of confining feedback modulation to a single isolated module, GDF-Net integrates continuous spatial attention gates across all decoder skip connections, modulating intermediate features via continuous spatial priors: $\mathbf{F}_\ell \odot (1 + \mathcal{A}_\ell(\mathbf{m}_{\text{guide}}))$.
  \item \textbf{Universal Benchmark Validation Across 7 Modalities:} A rigorous, reproducible experimental protocol validating GDF-Net across 7 diverse biomedical datasets (CVC-ClinicDB, EM-Dataset, Kvasir-SEG, ISIC-2018, CHASE-DB1, DRIVE transfer, and DSB-2018).
\end{enumerate}

\subsection{Summary of Contributions}
The contributions of this work are summarized as follows:
\begin{enumerate}
  \item \textbf{Theoretical Diagnosis of the Feedback Trap:} We provide the first systematic mathematical formulation and empirical demonstration of the Feedback Trap in recurrent vision networks, diagnosing internal attention gradient paralysis and backward recursive error amplification.
  \item \textbf{The Feedback Firewall \& GDF-Net Architecture:} We establish the Gradient Decoupling principle and propose GDF-Net, integrating a 4-channel decoupled guidance stem with a Multi-scale Gradient-Decoupled Feedback Decoder to enable stable, multi-resolution feature modulation.
  \item \textbf{Comprehensive 7-Dataset Benchmark Superiority:} We evaluate GDF-Net on 7 benchmark datasets across 4 distinct biomedical imaging domains. GDF-Net outperforms standard baseline architectures across 100\% of evaluated benchmarks, achieving \textbf{0.9427} Dice on CVC-ClinicDB (+4.69\,pp), \textbf{0.9518} on EM-Dataset (+17.02\,pp), \textbf{0.8880} on Kvasir-SEG (+0.28\,pp), \textbf{0.9000} on ISIC-2018 (+2.69\,pp), \textbf{0.8953} on DSB-2018 (+2.32\,pp), \textbf{0.7932} on DRIVE (+3.14\,pp), and slashing false positive rates by \textbf{65\%} on CHASE-DB1.
  \item \textbf{Iterative Dynamics and Reproducibility:} We analyze the convergence dynamics across feedback iterations, demonstrating that true iterative boundary refinement is unlocked only under gradient decoupling, and provide verified, reproducible open-source implementations for all benchmarks.
\end{enumerate}
